% figures/fig2_rama_a.tex
\begin{tikzpicture}[
    >=Stealth,
    % Estilo idéntico al bloque principal de la Figura 1
    block/.style={
        rectangle,
        draw=black!70,
        fill=black!2,
        rounded corners=4pt,
        line width=0.8pt,
        text width=6.2cm,
        align=center,
        inner sep=8pt,
        minimum height=1.35cm,
        font=\small\sffamily
    },
    % Estilo idéntico a las cajas terminales de las ramas de la Figura 1
    branchbox/.style={
        rectangle,
        draw=black!70,
        fill=black!4,
        rounded corners=4pt,
        line width=0.9pt,
        text width=4.6cm,
        align=center,
        inner sep=7pt,
        font=\small\sffamily
    },
    % Flechas principales
    arrow/.style={
        ->,
        line width=1.05pt,
        draw=black!75,
        shorten >= 2pt,
        shorten <= 2pt
    },
    % Líneas de unión sin punta
    line/.style={
        line width=1.05pt,
        draw=black!75
    },
    % Operador circular idéntico al de la Figura 1
    oplabel/.style={
        circle,
        draw=black!85,
        fill=white,
        line width=0.9pt,
        inner sep=3.5pt,
        font=\normalsize\bfseries
    }
]

    % =========================================================================
    % Nivel 0: Entrada de la Rama A (y = 0.0)
    % =========================================================================
    \node[block] (inA) at (0, 0) {
        \textbf{Vector de estado $\vec{P}$}\\
        \footnotesize Entrada a la Rama A (Banda estrecha en $L^2$)
    };

    % Operador de Fase Cero (y = -1.8)
    \node[oplabel] (opF) at (0, -1.8) {$\hat{\mathcal{F}}_{\text{ZPh}}$};

    % =========================================================================
    % Nivel 1: Señal Filtrada (y = -3.6)
    % =========================================================================
    \node[block] (filtered) at (0, -3.6) {
        \textbf{Señal filtrada de fase nula}\\
        \footnotesize $\mathbf{y}(t) = \mathbf{y}(t) \pm \Delta\mathbf{y}$ ($\tau_g \equiv 0$)
    };

    % Conexiones superiores
    \draw[arrow] (inA) -- (opF);
    \draw[arrow] (opF) -- (filtered);

    % =========================================================================
    % TRIFURCACIÓN SIMÉTRICA (3 VÍAS)
    % =========================================================================
    % 1. Tronco vertical descendente desde la señal filtrada
    \coordinate (split) at (0, -4.8);
    \draw[line] (filtered.south) -- (split);

    % 2. Barra horizontal de distribución (de -5.4 a +5.4)
    \coordinate (leftbar) at (-5.4, -4.8);
    \coordinate (rightbar) at (5.4, -4.8);
    \draw[line] (leftbar) -- (rightbar);

    % 3. Bajadas verticales hacia las 3 vías (x = -5.4, x = 0.0, x = +5.4)
    \draw[arrow] (-5.4, -4.8) -- (-5.4, -5.6);
    \draw[arrow] (0.0, -4.8) -- (0.0, -5.6);
    \draw[arrow] (5.4, -4.8) -- (5.4, -5.6);

    % =========================================================================
    % Nivel 2: Las 3 Vías de la Rama A (y = -7.2)
    % =========================================================================
    \node[branchbox] (via1) at (-5.4, -7.2) {
        \textbf{Vía 1: Cinemática Doppler}\\
        \footnotesize Transformada de Hilbert $\hat{\mathcal{H}}$\\
        Vector analítico $\mathbf{z}(t) \in \mathbb{C}^N$\\
        Derivada regularizada $\hat{\mathcal{D}}_{\Phi, SG}$
    };

    \node[branchbox] (via2) at (0.0, -7.2) {
        \textbf{Vía 2: Régimen Cuasi-Anecoico}\\
        \footnotesize Enventanado causal $\hat{T}_{\text{gate}}$ ($K=3$)\\
        Demodulación Lock-In $\hat{\mathcal{P}}_{f_0}$\\
        Integral de energía $\hat{\mathcal{E}}$
    };

    \node[branchbox] (via3) at (5.4, -7.2) {
        \textbf{Vía 3: Detección ToA}\\
        \footnotesize Ventana guarda $\hat{\mathcal{G}}_{\text{blank}}$\\
        Discriminador pendiente $\hat{\mathcal{C}}_{zc}$\\
        Interpolador secante $\hat{\mathcal{I}}_{zc}$
    };

\end{tikzpicture}